% GENERATED FILE. Edit canonical lessons, then run npm run generate:books.
% canonical-source-hash: fnv1a64:9a64bb122d9d239b
% canonical-lessons: ML-C297-maida, ML-C297-aval, ML-C297-kun, ML-C297-peraykka, ML-C297-nellikka

\chapter{Refined white flour, Flattened rice, Mushroom, Guava, Indian gooseberry}
\label{ch:ml-refined-white-flour-flattened-rice-mushroom-guava-indian-gooseberry}
\hlchaptermodality{297}

\begin{chapteropening}
\textbf{By the end of this chapter:} \emph{I can say refined white flour, flattened rice, a mushroom, a guava and an Indian gooseberry.}

Complete the last lesson of chapter 297: I can say refined white flour, flattened rice, a mushroom, a guava and an Indian gooseberry.
\end{chapteropening}

\section[\texorpdfstring{\ml{മൈദ} (maida)}{മൈദ (maida)}]{\ml{മൈദ} (maida) — refined white flour}

\label{lesson:ML-C297-maida}



\begin{quote}
Before the new one: say the Malayalam for a carpet, then the Malayalam for a bedsheet.
\end{quote}

\subsection*{You'll want to know: \ml{മൈദ}}
\textbf{\ml{മൈദ}} — \emph{maida} — \textquotedblleft{}refined white flour\textquotedblright{}.

\begin{grammarlens}[title={What you've built}]
One more word for this chapter.
\end{grammarlens}

\subsection*{Your turn}
\begin{itemize}
\raggedright
  \item \emph{Say it:} \emph{maida}
  \item \emph{Say it:} \emph{maida}, once more
  \item \emph{Say it:} say \emph{maida}
  \item \emph{Recall:} say the Malayalam for a carpet, then the Malayalam for a bedsheet, then say \emph{maida} again
\end{itemize}

\begin{tcolorbox}[breakable,colback=teal!4,colframe=teal!35!black,title={Before you move on}]
What does \ml{മൈദ} mean? (\textquotedblleft{}Refined white flour\textquotedblright{}.) Say it once more.
\end{tcolorbox}
\section[\texorpdfstring{\ml{അവൽ} (aval)}{അവൽ (aval)}]{\ml{അവൽ} (aval) — flattened rice}

\label{lesson:ML-C297-aval}



\begin{quote}
Before the new one: say the Malayalam for a bedsheet, then the Malayalam for refined white flour.
\end{quote}

\subsection*{You'll want to know: \ml{അവൽ}}
\textbf{\ml{അവൽ}} — \emph{aval} — \textquotedblleft{}flattened rice\textquotedblright{}.

\begin{grammarlens}[title={What you've built}]
One more word for this chapter.
\end{grammarlens}

\subsection*{Your turn}
\begin{itemize}
\raggedright
  \item \emph{Say it:} \emph{aval}
  \item \emph{Say it:} \emph{aval}, once more
  \item \emph{Say it:} say \emph{aval}
  \item \emph{Recall:} say the Malayalam for a bedsheet, then the Malayalam for refined white flour, then say \emph{aval} again
\end{itemize}

\begin{tcolorbox}[breakable,colback=teal!4,colframe=teal!35!black,title={Before you move on}]
What does \ml{അവൽ} mean? (\textquotedblleft{}Flattened rice\textquotedblright{}.) Say it once more.
\end{tcolorbox}
\section[\texorpdfstring{\ml{കൂൺ} (kūṇ)}{കൂൺ (kūṇ)}]{\ml{കൂൺ} (kūṇ) — a mushroom}

\label{lesson:ML-C297-kun}



\begin{quote}
Before the new one: say the Malayalam for refined white flour, then the Malayalam for flattened rice.
\end{quote}

\subsection*{You'll want to know: \ml{കൂൺ}}
\textbf{\ml{കൂൺ}} — \emph{kūṇ} — \textquotedblleft{}a mushroom\textquotedblright{}.

\begin{grammarlens}[title={What you've built}]
One more word for this chapter.
\end{grammarlens}

\subsection*{Your turn}
\begin{itemize}
\raggedright
  \item \emph{Say it:} \emph{kūṇ}
  \item \emph{Say it:} \emph{kūṇ}, once more
  \item \emph{Say it:} say \emph{kūṇ}
  \item \emph{Recall:} say the Malayalam for refined white flour, then the Malayalam for flattened rice, then say \emph{kūṇ} again
\end{itemize}

\begin{tcolorbox}[breakable,colback=teal!4,colframe=teal!35!black,title={Before you move on}]
What does \ml{കൂൺ} mean? (\textquotedblleft{}A mushroom\textquotedblright{}.) Say it once more.
\end{tcolorbox}
\section[\texorpdfstring{\ml{പേരയ്ക്ക} (pēraykka)}{പേരയ്ക്ക (pēraykka)}]{\ml{പേരയ്ക്ക} (pēraykka) — a guava}

\label{lesson:ML-C297-peraykka}



\begin{quote}
Before the new one: say the Malayalam for flattened rice, then the Malayalam for a mushroom.
\end{quote}

\subsection*{You'll want to know: \ml{പേരയ്ക്ക}}
\textbf{\ml{പേരയ്ക്ക}} — \emph{pēraykka} — \textquotedblleft{}a guava\textquotedblright{}.

\begin{grammarlens}[title={What you've built}]
One more word for this chapter.
\end{grammarlens}

\subsection*{Your turn}
\begin{itemize}
\raggedright
  \item \emph{Say it:} \emph{pēraykka}
  \item \emph{Say it:} \emph{pēraykka}, once more
  \item \emph{Say it:} say \emph{pēraykka}
  \item \emph{Recall:} say the Malayalam for flattened rice, then the Malayalam for a mushroom, then say \emph{pēraykka} again
\end{itemize}

\begin{tcolorbox}[breakable,colback=teal!4,colframe=teal!35!black,title={Before you move on}]
What does \ml{പേരയ്ക്ക} mean? (\textquotedblleft{}A guava\textquotedblright{}.) Say it once more.
\end{tcolorbox}
\section[\texorpdfstring{\ml{നെല്ലിക്ക} (nellikka)}{നെല്ലിക്ക (nellikka)}]{\ml{നെല്ലിക്ക} (nellikka) — an Indian gooseberry}

\label{lesson:ML-C297-nellikka}



\begin{quote}
Before the new one: say the Malayalam for a mushroom, then the Malayalam for a guava.
\end{quote}

\subsection*{You'll want to know: \ml{നെല്ലിക്ക}}
\textbf{\ml{നെല്ലിക്ക}} — \emph{nellikka} — \textquotedblleft{}an Indian gooseberry\textquotedblright{}.

\begin{grammarlens}[title={What you've built}]
One more word for this chapter.
\end{grammarlens}

\subsection*{Your turn}
\begin{itemize}
\raggedright
  \item \emph{Say it:} \emph{nellikka}
  \item \emph{Say it:} \emph{nellikka}, once more
  \item \emph{Say it:} say \emph{nellikka}
  \item \emph{Recall:} say the Malayalam for a mushroom, then the Malayalam for a guava, then say \emph{nellikka} again
\end{itemize}

\begin{tcolorbox}[breakable,colback=teal!4,colframe=teal!35!black,title={Before you move on}]
What does \ml{നെല്ലിക്ക} mean? (\textquotedblleft{}An Indian gooseberry\textquotedblright{}.) Say it once more.
\end{tcolorbox}
